\RequirePackage[l2tabu, orthodox]{nag}
\RequirePackage{silence}
\WarningFilter{nag}{There is no environment ``centering'' }%nag complains because beamer titlepage uses a centering environment
\WarningFilter{nag}{1 complaints in total}
\WarningFilter{ifpdf}{Someone has redefined \pdfoutput}
\WarningFilter{fmtcount}{\ordinal already defined use \FCordinal instead}
\documentclass[english, french]{beamer}
%INSTALL

%avoids a warning
\usepackage[log-declarations=false]{xparse}
\usepackage{fontspec} %font selecting commands
\usepackage{xunicode}
%warn about missing characters
\tracinglostchars=2

%REDAC
\usepackage{booktabs}
\usepackage{calc}

\usepackage{mathtools} %load this before babel!
	\mathtoolsset{showonlyrefs,showmanualtags}

\usepackage{babel}
%suppresses the warning about frenchb not modifying the captions (“—” to “:” in “Figure 1 – Legend”).
	\frenchbsetup{AutoSpacePunctuation=false,SuppressWarning=false}

\usepackage[super]{nth}
\usepackage{listings} %typeset source code listings
	\lstset{language=XML,tabsize=2,literate={"}{{\tt"}}1,captionpos=b}
\usepackage[nolist,smaller,printonlyused]{acronym}%,smaller option produces warnings from relsize in some cases, it seems.
\usepackage[nodayofweek]{datetime}%must be loaded after the babel package
\usepackage{xspace}
\usepackage{hyperref}% option pdfusetitle must be introduced here, not in hypersetup.
%breaklinks makes links on multiple lines into different PDF links to the same target.
%colorlinks (false): Colors the text of links and anchors. The colors chosen depend on the the type of link. In spite of colored boxes, the colored text remains when printing.
%linkcolor=black: this leaves other links in colors, e.g. refs in green, don't print well.
%pdfborder (0 0 1, set to 0 0 0 if colorlinks): width of PDF link border
%hidelinks
\hypersetup{breaklinks,bookmarksopen,colorlinks=true,urlcolor=blue,linkcolor=,hyperfigures=true}
% hyperref doc says: Package bookmark replaces hyperref’s bookmark organization by a new algorithm (...) Therefore I recommend using this package.
\usepackage{bookmark}

% center floats by default, but do not use with float
% \usepackage{floatrow}
% \makeatletter
% \g@addto@macro\@floatboxreset\centering
% \makeatother
\usepackage{ragged2e} %new com­mands \Cen­ter­ing, \RaggedLeft, and \RaggedRight and new en­vi­ron­ments Cen­ter, FlushLeft, and FlushRight, which set ragged text and are eas­ily con­fig­urable to al­low hy­phen­ation (the cor­re­spond­ing com­mands in LaTeX, all of whose names are lower-case, pre­vent hy­phen­ation al­to­gether). 
\usepackage{siunitx} %[expproduct=tighttimes, decimalsymbol=comma]
\sisetup{detect-all}% to detect e.g. when in math mode (use a math font)
\usepackage{braket} %for \Set
\usepackage{natbib}

\usepackage{amsmath,amsthm}
% \usepackage{amsfonts} %not required?
% \usepackage{dsfont} %for what?
%unicode-math overwrites the following commands from the mathtools package: \dblcolon, \coloneqq, \Coloneqq, \eqqcolon. Using the other colon-like commands from mathtools will lead to inconsistencies. Plus, Using \overbracket and \underbracke from mathtools package. Use \Uoverbracket and \Uunderbracke for original unicode-math definition.
%use exclusively \mathbf and choose math bold style below.
\usepackage[warnings-off={mathtools-colon, mathtools-overbracket}, bold-style=ISO]{unicode-math}

\defaultfontfeatures{
	Fractions=On,
	Mapping=tex-text% to turn "--" into dashes, useful for bibtex%%
}
\defaultfontfeatures[\rmfamily, \sffamily]{
	Fractions=On,
	Mapping=% to leave " alone (disable the default mapping tex-text; requires loading the font afterwards?)
}
\newfontfamily\xitsfamily{XITS}
\newfontfamily\texgyretermesfamily{TeX Gyre Termes}
\newfontfamily\texgyreherosfamily{TeX Gyre Heros}
\newfontfamily\lmfamily{Latin Modern Roman}
\setmainfont{Latin Modern Roman}
\setsansfont{Latin Modern Sans}
\setmonofont{Latin Modern Mono}
\defaultfontfeatures{}%disable default font features to avoid warnings with math fonts.
\setmathfont{XITS Math}
\setmathfont[range={\mathcal,\mathbfcal},StylisticSet=1]{XITS Math}

\usepackage{cleveref}% cleveref should go "laster" than hyperref
%GRAPHICS
\usepackage{pgf}
\usepackage{pgfplots}
	\usetikzlibrary{matrix,fit,plotmarks,calc,trees,shapes.geometric,positioning,plothandlers}
\pgfplotsset{compat=1.11}
\usepackage{graphicx}

\graphicspath{{graphics/},{graphics-dm/}}
\DeclareGraphicsExtensions{.pdf}

%HACKING
\usepackage{printlen}
\uselengthunit{mm}
% 	\newlength{\templ}% or LenTemp?
% 	\setlength{\templ}{6 pt}
% 	\printlength{\templ}
\usepackage{etoolbox} %for addtocmd command
\usepackage{scrhack}% load at end. Corrects a bug in float package, which is outdated but might be used by other packages
\usepackage{xltxtra} %somebody said that this is loaded by fontspec, but does not seem correct: if not loaded explicitly, does not appear in the log and \showhyphens is not corrected.

%Beamer-specific
%ADD
\usepackage{appendixnumberbeamer}
%\setbeamersize{text margin left=0.1cm, text margin right=0.1cm} 
\setbeamertemplate{navigation symbols}{}
\usetheme{BrusselsBelgium}
\usefonttheme{professionalfonts}


\newcommand{\R}{ℝ}
\newcommand{\N}{ℕ}
\newcommand{\Z}{ℤ}
\newcommand{\card}[1]{\lvert{#1}\rvert}
\newcommand{\powerset}[1]{\mathscr{P}(#1)}%\mathscr rather than \mathcal: scr is rounder than cal (at least in XITS Math).
\newcommand{\suchthat}{\;\ifnum\currentgrouptype=16 \middle\fi|\;}
%\newcommand{\Rplus}{\reels^+\xspace}

\AtBeginDocument{%
	\renewcommand{\epsilon}{\varepsilon}
% we want straight form of \phi for mathematics, as recommended in UTR #25: Unicode support for mathematics.
%	\renewcommand{\phi}{\varphi}
}

% with amssymb, but I don’t want to use amssymb just for that.
% \newcommand{\restr}[2]{{#1}_{\restriction #2}}
%\newcommand{\restr}[2]{{#1\upharpoonright}_{#2}}
\newcommand{\restr}[2]{{#1|}_{#2}}%sometimes typed out incorrectly within \set.
%\newcommand{\restr}[2]{{#1}_{\vert #2}}%\vert errors when used within \Set and is typed out incorrectly within \set.
\DeclareMathOperator*{\argmax}{arg\,max}
\DeclareMathOperator*{\argmin}{arg\,min}


%ARG TH
\newcommand{\AF}{\mathcal{AF}}
\newcommand{\labelling}{\mathcal{L}}
\newcommand{\labin}{\textbf{in}\xspace}
\newcommand{\labout}{\textbf{out}}
\newcommand{\labund}{\textbf{undec}\xspace}
\newcommand{\nonemptyor}[2]{\ifthenelse{\equal{#1}{}}{#2}{#1}}
\newcommand{\gextlab}[2][]{
	\labelling{\mathcal{GE}}_{(#2, \nonemptyor{#1}{\ibeatsr{#2}})}
}
\newcommand{\allargs}{A^*}
\newcommand{\args}{A}
\newcommand{\ar}{a}

%MCDA+Arg
\newcommand{\dm}{d}
\newcommand{\ileadsto}{\leadsto}
\newcommand{\mleadsto}[1][\eta]{\leadsto_{#1}}
\newcommand{\ibeats}{\vartriangleright}
\newcommand{\mbeats}[1][\eta]{\vartriangleright_{#1}}


%MISC
\newcommand{\lequiv}{\Vvdash}
\newcommand{\weightst}{W^{\,t}}

%MCDA classical
\newcommand{\crits}{\mathcal{J}}
\newcommand{\altspace}{\mathbb{A}}
\newcommand{\alts}{A}

%Sorting
\newcommand{\cats}{\mathcal{C}}
\newcommand{\catssubsets}{2^\cats}
\newcommand{\catgg}{\vartriangleright}
\newcommand{\catll}{\vartriangleleft}
\newcommand{\catleq}{\trianglelefteq}
\newcommand{\catgeq}{\trianglerighteq}
\newcommand{\alttoc}[2][x]{(#1 \xrightarrow{} #2)}
\newcommand{\alttocat}[3]{(#2 \xrightarrow{#1} #3)}
\newcommand{\alttoI}{(x \xrightarrow{} \left[\underline{C_x}, \overline{C_x}\right])}
\newcommand{\alttocatdm}[3][t]{\left(#2 \thinspace \raisebox{-3pt}{$\xrightarrow{#1}$}\thinspace #3\right)}
\newcommand{\alttocatatleast}[2]{\left(#1 \thinspace \raisebox{-3pt}{$\xrightarrow[]{≥}$}\thinspace #2\right)}
\newcommand{\alttocatatmost}[2]{\left(#1 \thinspace \raisebox{-3pt}{$\xrightarrow[]{≤}$}\thinspace #2\right)}

\newcommand{\source}{\scriptsize}
\newcommand{\commentOC}[1]{{\selectlanguage{french}{\todo{OC : #1}}}}
%Or: \todo[color=green!40]

%this probably requires outdated float package, see doc KomaScript for an alternative.
% \newfloat{program}{t}{lop}
% \floatname{program}{PM}

%\crefname{axiom}{axiom}{axioms}%might be needed for workaround bug in cref when defining new theorems?

%\ifdefined\theorem\else
%\newtheorem{theorem}{\iflanguage{english}{Theorem}{Théorème}}
%\fi

%which line breaks are chosen: accept worse lines, therefore reducing risk of overfull lines. Default = 200
\tolerance=2000
%accept overfull hbox up to...
\hfuzz=2cm
%reduces verbosity about the bad line breaks
\hbadness 5000
%sloppy sets tolerance to 9999
\apptocmd{\sloppy}{\hbadness 10000\relax}{}{}

% WRITING
%\newcommand{\ie}{i.e.\@\xspace}%to try
%\newcommand{\eg}{e.g.\@\xspace}
%\newcommand{\etal}{et al.\@\xspace}
\newcommand{\ie}{i.e.\ }
\newcommand{\eg}{e.g.\ }
\newcommand{\mkkOK}{\checkmark}%\color{green}{\checkmark}
\newcommand{\mkkREQ}{\ding{53}}%requires pifont?%\color{green}{\checkmark}
\newcommand{\mkkNO}{}%\text{\color{red}{\textsf{X}}}

\makeatletter
\newcommand{\boldor}[2]{%
	\ifnum\strcmp{\f@series}{bx}=\z@
		#1%
	\else
		#2%
	\fi
}
\newcommand{\textstyleElProm}[1]{\boldor{\MakeUppercase{#1}}{\textsc{#1}}}
\makeatother
\newcommand{\electre}{\textstyleElProm{Électre}\xspace}
\newcommand{\electreIv}{\textstyleElProm{Électre Iv}\xspace}
\newcommand{\electreIV}{\textstyleElProm{Électre IV}\xspace}
\newcommand{\electreIII}{\textstyleElProm{Électre III}\xspace}
\newcommand{\electreTRI}{\textstyleElProm{Électre Tri}\xspace}
% \newcommand{\utadis}{\texorpdfstring{\textstyleElProm{utadis}\xspace}{UTADIS}}
% \newcommand{\utadisI}{\texorpdfstring{\textstyleElProm{utadis i}\xspace}{UTADIS I}}

%TODO
% \newcommand{\textstyleElProm}[1]{{\rmfamily\textsc{#1}}} 


\newlength{\GraphsNodeSep}
\setlength{\GraphsNodeSep}{7mm}

% MCDA Drawing Sorting
\newlength{\MCDSCatHeight}
\setlength{\MCDSCatHeight}{6mm}
\newlength{\MCDSAltHeight}
\setlength{\MCDSAltHeight}{4mm}
%separation between two vertical alts
\newlength{\MCDSAltSep}
\setlength{\MCDSAltSep}{2mm}
\newlength{\MCDSCatWidth}
\setlength{\MCDSCatWidth}{3cm}
\newlength{\MCDSEvalRowHeight}
\setlength{\MCDSEvalRowHeight}{6mm}
\newlength{\MCDSAltsToCatsSep}
\setlength{\MCDSAltsToCatsSep}{1.5cm}
\newcounter{MCDSNbAlts}
\newcounter{MCDSNbCats}
\newlength{\MCDSArrowDownOffset}
\setlength{\MCDSArrowDownOffset}{0mm}

\tikzset{/Graphs/dot/.style={
	shape=circle, fill=black, inner sep=0, minimum size=1mm
}}
\tikzset{/MC/D/S/alt/.style={
	shape=rectangle, draw=black, inner sep=0, minimum height=\MCDSAltHeight, minimum width=2.5cm, anchor=north east
}}
\tikzset{MC/D/S/pref/.style={
	shape=ellipse, draw=gray, thick
}}
\tikzset{/MC/D/S/cat/.style={
	shape=rectangle, draw=black, inner sep=0, minimum height=\MCDSCatHeight, minimum width=\MCDSCatWidth, anchor=north west
}}
\tikzset{/MC/D/S/evals matrix/.style={
	matrix, row sep=-\pgflinewidth, column sep=-\pgflinewidth, nodes={shape=rectangle, draw=black, inner sep=0mm, text depth=0.5ex, text height=1em, minimum height=\MCDSEvalRowHeight, minimum width=12mm}, nodes in empty cells, matrix of nodes, inner sep=0mm, outer sep=0mm, row 1/.style={nodes={draw=none, minimum height=0em, text height=, inner ysep=1mm}}
}}

% Beliefs
\tikzset{/Beliefs/D/S/attacker/.style={
	shape=rectangle, draw, minimum size=8mm
}}
\tikzset{/Beliefs/D/S/supporter/.style={
	shape=circle, draw
}}

\newcommand{\tikzmark}[1]{%
	\tikz[overlay, remember picture, baseline=(#1.base)] \node (#1) {};%
}


\begin{acronym}
\acro{AMCD}{Aide Multicritère à la Décision}
\acro{ASA}{Argument Strength Assessment}
\acro{DA}{Decision Analysis}
\acro{DM}{Decision Maker}
\acro{DPr}{Deliberated Preferences}
\acro{DRSA}{Dominance-based Rough Set Approach}
\acro{DSS}{Decision Support Systems}
\acrodefplural{DSS}{Decision Support Systems}
% \newacroplural{DSS}[DSSes]{Decision Support Systems}
\acro{EJOR}{European Journal of Operational Research}
\acro{LNCS}{Lecture Notes in Computer Science}
\acro{MCDA}{Multicriteria Decision Aid}
\acro{MIP}{Mixed Integer Program}
\acro{NCSM}{Non Compensatory Sorting Model}
\acro{PL}{Programme Linéaire}
\acro{PLNE}{Programme Linéaire en Nombres Entiers}
\acro{PM}{Programme Mathématique}
\acro{MP}{Mathematical Program}
\acro{MIP}{Mixed Integer Program}
% \newacroplural{PM}{Programmes Mathématiques}
%acrodefplural since version 1.35, my debian has \ProvidesPackage{acronym}[2009/01/25, v1.34, Support for acronyms (Tobias Oetiker)]
\acrodefplural{PM}{Programmes Mathématiques}
\acro{PMML}{Predictive Model Markup Language}
\acro{RESS}{Reliability Engineering \& System Safety}
\acro{SMAA}{Stochastic Multicriteria Acceptability Analysis}
\acro{URPDM}{Uncertainty and Robustness in Planning and Decision Making}
\acro{XML}{Extensible Markup Language}
\end{acronym}


\title{Conception d’applications internet}
\subtitle{Services Web REST}
\subject{Java EE}
\keywords{REST, JAXB, Web services}
\author{Olivier Cailloux}
\institute[LAMSADE]{LAMSADE, Université Paris-Dauphine}
\date{}

\begin{document}
\bibliographystyle{apalike}

\begin{frame}[plain]
	\tikz[remember picture,overlay]{
		\path (current page.south west) node[anchor=south west, inner sep=0] {
			\includegraphics[height=1cm]{LAMSADE95.jpg}
		};
		\path (current page.south) ++ (0, 1mm) node[anchor=south, inner sep=0] {
			\includegraphics[height=9mm]{Dauphine.jpg}
		};
		\path (current page.south east) node[anchor=south east, inner sep=0] {
			\includegraphics[height=1cm]{PSL.png}
		};
	}
   \titlepage
\end{frame}
\addtocounter{framenumber}{-1}

\section{JAX-RS}
\subsection{Introduction}
\begin{frame}
	\frametitle{Services Web}
	\begin{itemize}
		\item SOAP (Simple Object Access Protocol) : standard W3C
		\item Définit SOAP enveloppe pour transmission (basée sur XML), architecture de liaison avec protocoles sous-jacents, liaison HTTP, compression de messages…
		\item REST (Representational State Transfer) : contraintes architecturales
	\end{itemize}
	Cf. \href{http://www.w3.org/2005/Talks/1115-hh-k-ecows/}{présentation} : Reconciling Web Services and REST Services
\end{frame}

\begin{frame}
	\frametitle{Méthodes HTTP}
	\begin{itemize}
		\item Basé sur \emph{ressource} (identifiée par URI)
	\end{itemize}
	\begin{block}{Exemples de méthodes}
		\begin{itemize}
			\item GET : demander représentation d’une ressource {\tiny ou une partie}
			\item POST : traiter représentation fournie selon ressource indiquée
			\item PUT : remplacer la ressource par la représentation indiquée
			\item DELETE : effacer le lien entre la ressource et sa fonction
		\end{itemize}
	\end{block}
	\begin{itemize}
		\item Idempotent ? \pause Répétition produit même effet (GET, PUT, DELETE)
		\item Safe ? \pause Pas de changement d’état (GET)
		\item Représentation sous formes diverses (\href{http://www.iana.org/assignments/media-types/}{liste} \href{http://tools.ietf.org/html/rfc7231\#section-3.1.1.1}{media-types})
	\end{itemize}
\end{frame}

\begin{frame}
	\frametitle{REST en Java EE : aperçu}
	\begin{itemize}
		\item JAX-RS
		\item Sorte de servlet plus haut-niveau
		\item Classe ressource : l’annoter \jeeref[@]{javax.ws.rs/Path}
		\item Méthode ressource :  l’annoter \jeeref[@]{javax.ws.rs/GET}, etc.
		\item \texttt{@Path} sur méthode : est une sub-resource method {\tiny ou sub-resource locator}
	\end{itemize}
\end{frame}

\subsection{Implémentation}
\begin{frame}
	\frametitle{Méthodes ressource}
	\begin{itemize}
		\item Méthode renvoyant une réponse : annoter \jeeref[@]{javax.ws.rs/Produces}, tableau de \texttt{String} (cf. \jeeref{javax.ws.rs.core/MediaType})
		\item Choix méthode appelée par le conteneur ? \pause selon HTTP \texttt{accept} et annotations \texttt{@Produces} \pause
		\item Type retour méthode ? \pause Doit être convertible en type mime adéquat \pause
		\item Conversion par \emph{Entity provider} \og{}writer\fg{} embarqué ou maison
		\item Ou : renvoyer \jeeref{javax.ws.rs.core/Response} {\tiny ou \jeeref{javax.ws.rs.core/GenericEntity}}
		\item Statut ? \pause Mis automatiquement si pas mis {\tiny et pas d’exception} : \texttt{null} ⇒ 204 No Content ; sinon 200 OK
	\end{itemize}
\end{frame}

\begin{frame}
	\frametitle{Entity providers}
	\begin{itemize}
		\item Provider : classe annotée \jeeref[@]{javax.ws.rs.ext/Provider} et {\tiny éventuellement} \jeeref[@]{javax.ws.rs/Produces}
		\item Entity provider \og{}writer\fg{} : provider qui implémente \jeeref{javax.ws.rs.ext/MessageBodyWriter}\texttt{<MyType>}
		\item Entity provider fournis, exemple \texttt{text/plain} ⇔ \texttt{String} ; \texttt{MultivaluedMap<String, String>} ⇔ \texttt{application/x-www-form-urlencoded}…
		\item{} {\tiny [text/xml , application/xml, application/*+xml] ⇔ [javax.xml.transform.Source, javax.xml.bind.JAXBElement, application-supplied JAXB classes]} %TODO
		\item Méthode consommant un contenu : annoter \jeeref[@]{javax.ws.rs/Consumes}, tableau de \texttt{String} (cf. \jeeref{javax.ws.rs.core/MediaType})
		\item Lecture du contenu ? \pause Méthode a un paramètre non annoté de type convertible en type mime adéquat \pause
		\item Entity provider \og{}reader\fg{} : classe annotée \jeeref[@]{javax.ws.rs.ext/Provider} qui implémente \jeeref{javax.ws.rs.ext/MessageBodyReader}
	\end{itemize}
\end{frame}

\subsection{Paramètres}
\begin{frame}
	\frametitle{Paramètres}
	Accès aux paramètres autres que contenu requête HTTP :
	\begin{itemize}
		\item \texttt{\jeeref[@]{javax.ws.rs/DefaultValue}("me") \jeeref[@]{javax.ws.rs/QueryParam}("name") String na}
		\item Ou classe avec paramètre constructeur \texttt{String} {\tiny ou collections}
		\item Convertisseurs de base fournis : \texttt{@QueryParam("zevalue") int value}, …
		\item Penser au défaut explicite ! (Sinon défaut pour type)
		\item Paramètre dans l’URI : \texttt{@Path("users/\{username\}")} {\tiny ou \texttt{@Path("users/\{username: regexp\}")}, par défaut [\^{}/]+?}
		\item \texttt{\jeeref[@]{javax.ws.rs/PathParam}("username") String usernameVar;}
		\item \jeeref[@]{javax.ws.rs/HeaderParam}, \jeeref[@]{javax.ws.rs/CookieParam}, \jeeref[@]{javax.ws.rs/MatrixParam}, \jeeref[@]{javax.ws.rs/FormParam}
		\item \jeeref[@]{javax.ws.rs.core/Context} \jeeref{javax.ws.rs.core/UriInfo} \texttt{ui} ; \texttt{@Context \jeeref{javax.ws.rs.core/HttpHeaders} hh}
		\item Aussi sur champs {\tiny ou constructeur pour certains} (restriction ? \pause pour classes ressources de scope requête)
	\end{itemize}
\end{frame}


\subsection{Cycles de vie}
\begin{frame}
	\frametitle{Cycles de vie}
	\begin{itemize}
		\item Scopes CDI permis {\tiny en fonction de l’implémentation JAX-RS} {\tiny spécification ambigüe concernant Provider mais fonctionne en pratique}
		\item Recommandé : toujours préciser scope sur classes ressources et providers
		\item Permet protection contre accès concurrents
	\end{itemize}
\end{frame}

\subsection{Exception}
\begin{frame}
	\frametitle{Exceptions}
	\begin{itemize}
		\item Provider maison : implémenter \jeeref{javax.ws.rs.ext/ExceptionMapper}\texttt{<MyException>} et renvoyer une Réponse
		\item Ou lancer \jeeref{javax.ws.rs/WebApplicationException}, indiquer réponse
		\item \texttt{WebApplicationException} enveloppée dans \jeeref{javax.servlet/ServletException} pour la propager {\tiny pour implémentation appuyée sur servlets}
	\end{itemize}
\end{frame}

\subsection{Packaging}
\begin{frame}
	\frametitle{Packaging}
	\begin{itemize}
		\item \jeeref[@]{javax.ws.rs/ApplicationPath} sur sous-classe de \jeeref{javax.ws.rs.core/Application} (capture value/*)
		\item Dans \texttt{.war}
		\item Application sous-classe : singleton {\tiny (@ApplicationScope obligatoirement)}
	\end{itemize}
\end{frame}

\subsection{Voir aussi}
\begin{frame}
	\frametitle{Voir aussi}
	\begin{itemize}
		\item Validation ; Asynchronicité ; Injection de contexte% (dans spec)
		\item \jeeref{javax.ws.rs.core/UriBuilder} peut être utile
		\item Considérer implémenter HEAD (sinon délégation à GET)
%		\item ParamConverter is available via a registered ParamConverterProvider, pour champ @TrucParam MyType param;
		\item \href{http://www.w3.org/Submission/wadl/}{WADL}, description formelle de services REST
		\item Une classe ressource peut être un Session bean (quelle sorte ? \pause SLSB ou Singleton : gestion cycle de vie par le conteneur)
	\end{itemize}
\end{frame}

\subsection{Client}
\begin{frame}[fragile]
	\frametitle{Client}
	\begin{itemize}
		\item \jeeref{javax.ws.rs.client/Client} : point d’entrée
		\item \jeeref{javax.ws.rs.client/WebTarget} pour manipuler des URIs
	\end{itemize}
	\begin{block}{Exemple}%
%		\lstinputlisting[language=Java,basicstyle=\NoAutoSpacing]{Client.java}
		\begin{lstlisting}[language=Java, basicstyle=\small\NoAutoSpacing, aboveskip=0pt, belowskip=0pt]
Client client = ClientBuilder.newClient();
WebTarget t1 = client.target("http://.../users");
WebTarget t2 = t1.path("{uId}");
WebTarget t3 = t2.resolveTemplate("uId", "123");
WebTarget t4 = t3.queryParam("prof", "1");
/** http://.../users/123?prof=1 */
String result = t4.request(MediaType.TEXT_PLAIN)
  .get(String.class);
client.close();
		\end{lstlisting}
	\end{block}
\end{frame}

\subsection{Références}
\begin{frame}
	\frametitle{Références}
	\begin{itemize}
		\item \href{http://dx.doi.org/10.1145/337180.337228}{Article} REST : Fielding \& Taylor, Principled Design of the Modern Web Architecture (\href{https://www.ics.uci.edu/~fielding/pubs/webarch_icse2000.pdf}{direct})%http://dx.doi.org.gate6.inist.fr/10.1145/337180.337228
		\item \href{http://www.w3.org/Protocols/}{RFC 7231} : HTTP/1.1, part 2: Semantics and Content
		\item The Java EE Tutorial: \href{http://docs.oracle.com/javaee/7/tutorial/partwebsvcs.htm}{Web Services}
		\item \href{https://jcp.org/en/jsr/detail?id=346}{JSR 346} (Context and Dependency Injection 1.1 et 1.2) (\href{http://download.oracle.com/otn-pub/jcp/cdi-1_2-mrel-eval-spec/cdi-1.2.pdf}{direct}).
		\item \href{https://jcp.org/en/jsr/detail?id=330}{JSR 330} (Dependency Injection) (\href{http://download.oracle.com/otn-pub/jcp/dependency_injection-1.0-final-oth-JSpec/dependency_injection-1_0-final-spec.zip}{direct}) : simplement \href{http://docs.oracle.com/javaee/7/api/javax/inject/Inject.html}{\texttt{@Inject}} et cie
		\item \href{https://jcp.org/en/jsr/detail?id=339} JSR 339 (JAX-RS 2.0) (\href{http://download.oracle.com/otn-pub/jcp/jaxrs-2_0_rev_A-mrel-eval-spec/jsr339-jaxrs-2.0-final-spec.pdf}{direct})
		\item \href{https://jcp.org/en/jsr/detail?id=342}{JSR 342} (Java EE 7) (\href{http://download.oracle.com/otn-pub/jcp/java_ee-7-mrel-eval-spec/JavaEE_Platform_Spec.pdf}{direct})
		\item \href{https://jcp.org/en/jsr/detail?id=345}{JSR 345} (EJB 3.2) (\href{http://download.oracle.com/otn-pub/jcp/ejb-3_2-fr-eval-spec/ejb-3_2-core-fr-spec.pdf}{direct})
	\end{itemize}
\end{frame}

\section{Exercices}
\begin{frame}[allowframebreaks]
	\frametitle{Exercices}
	\begin{itemize}
		\item Implémenter un GET qui renvoie une description d’un objet métier quelconque de votre projet (ou un objet jouet) comme type \texttt{text/plain}. La méthode doit renvoyer un \texttt{String}. Tester.
		\item Simuler un traitement long dans la méthode ressource (Thread sleep). Avec deux navigateurs différents, envoyer deux GET simultanés. Vérifier qu’elles sont bien traitées en parallèle (la deuxième n’attend pas la fin de la première avant de s’exécuter).
		\item Modifier la méthode ressource pour qu’elle renvoie l’objet lui-même. Implémenter un provider qui fournit sa description.
		\item Observer si les requêtes sont traitées en parallèle dans votre provider (simuler un traitement long dans le provider, envoyer deux GET simultanés avec deux navigateurs différents). Modifier la portée CDI du provider et observer l’impact.
		\item Insérer un état dans le provider. Le provider possède maintenant un compteur \texttt{i}. Il incrémente \texttt{i} en entrée d’appel, puis il indique la valeur de \texttt{i} dans la description renvoyée, puis il décrémente \texttt{i}. Le provider renverra-t-il toujours la valeur d’initialisation de \texttt{i} + 1, lors d’appels parallèles de clients web différents ? Prédire son comportement en fonction de la portée CDI qui vous lui affectez. Réfléchir aux avantages de différents choix.
	\end{itemize}
\end{frame}

\section{JAXB}
\begin{frame}
	\frametitle{Objectifs}
	\begin{itemize}
		\item Java Architecture for XML Binding
		\item Schema XML ⇔ Classes Java (⇒ xjc ; ⇐ schemagen)
		\item Instances XML ⇔ Instances Java (JAXB binding runtime)
		\item JAXB object : annotés
		\item JAX-RS fournit entity providers pour objets JAXB
%		\item http://www.xmlvalidation.com/test.xsd
	\end{itemize}
%	\begin{block}{}
%		ContextResolver<JAXBContext> annotated with @Provider. Context providers MAY return null from the getContext method if they do not wish to provide their context for a particular Java type. E.g. a JAXB context provider may wish to only provide the context for certain JAXB classes.%TODO 
%	\end{block}
\end{frame}

\begin{frame}
	\frametitle{Références JAXB}
	\begin{block}{Références}
		\begin{itemize}
			\item \href{https://docs.oracle.com/javase/tutorial/jaxb/index.html}{Tutoriel Oracle}
			\item \href{http://docs.oracle.com/javase/8/docs/technotes/guides/xml/jaxb/}{Docs Oracle}
		\end{itemize}
	\end{block}
	Dans eclipse :
	\begin{itemize}
		\item Désactiver le validateur est probablement nécessaire
		\item[⇒] cf. \href{https://bugs.eclipse.org/bugs/show_bug.cgi?id=437635}{bug} eclipse%ça fonctionne avec eclipse mars malgré le bug, en tous cas après désactivation du validateur
	\end{itemize}
\end{frame}

\appendix
\AtBeginSection{
}
\section{Licence}
\begin{frame}
	\frametitle{Licence}
	Cette présentation, et le code LaTeX associé, sont sous \href{http://opensource.org/licenses/MIT}{licence MIT}. Vous êtes libres de réutiliser des éléments de cette présentation, sous réserve de citer l’auteur.
	
	Le travail réutilisé est à attribuer à \href{http://www.lamsade.dauphine.fr/~ocailloux/}{Olivier Cailloux}, Université Paris-Dauphine.
	
	\small{(Ceci ne couvre pas les images incluses dans ce document, puisque je n’en suis généralement pas l’auteur.)}
\end{frame}
\end{document}
\begin{frame}
	\frametitle{}
	\begin{itemize}
		\item 
	\end{itemize}
	\begin{block}{}
		
	\end{block}
\end{frame}

\section{Bibliographie}
\begin{frame}[allowframebreaks]
	\frametitle{Bibliographie}
	\def\newblock{\hskip .11em plus .33em minus .07em}
% 	\bibliography{zotero}
\end{frame}

\section{Autres}
\begin{frame}
	\frametitle{}
	\begin{itemize}
		\item 
	\end{itemize}
	\begin{block}{}
		
	\end{block}
\end{frame}
\end{document}
